\chapter{Roadmap and known issues}\label{ch:roadmap}
\chapterstatus{complete}{2026-11-16}{product owner; \src{infra/aws}}

Planned improvements, in build order. Each item has an ID so commits and tests can refer to it. When an item ships, remove it from here and update \cref{ch:journeys,ch:rules}.

Priority: \textbf{P1} a user can bypass a rule, get stuck, or end up with wrong data · \textbf{P2} a real task is missing or awkward · \textbf{P3} polish.

Finished work is removed from this chapter. What shipped is described in \cref{ch:journeys,ch:rules}, and git history has the details (phases 1--4 merged on 2026-09-25).

\section{Next: check with users}

Run short task-based sessions with the seed accounts: book, reschedule, cancel then rebook, run a demo day, mark and submit, review. Record task success, errors and time taken. Turn what you find into new roadmap items.

After launch, \route{/platform/insights} and \route{/platform/feedback} give the same evidence continuously (\cref{sec:rules-usage}). Start with these questions, and change the event catalog as they are answered:

\begin{itemize}
  \item \textbf{Q1} Can students book without help? Booking funnel by device, the post-booking survey, booking-change requests.
  \item \textbf{Q2} Do reschedule and cancel confuse people? Reschedules vs cancellations rebooked within 30 minutes, the reschedule funnel and survey.
  \item \textbf{Q3} Is it fast enough on phones? Web Vitals by device; the slowest pages.
  \item \textbf{Q4} Do staff manage demo days and marking smoothly? Demo-to-marks time, unlocks, the marking funnel and survey.
  \item \textbf{Q5} Which features does nobody use? Pages and features by role.
  \item \textbf{Q6} Where does it break? Errors by page and release.
\end{itemize}

Known follow-up: enforce the content security policy (zod \code{jitless}, script nonces).

\section{Proposed: notifications and calendar (needs a decision)}

Feedback (2026-10-02): a subscribed feed is awkward, but per-demo calendar entries plus emails plus push could flood people with the same reminder. Proposal, one channel per purpose:

{\tablefont
\begin{xltabular}{\linewidth}{@{}P{30mm} P{16mm} P{30mm} P{30mm} L@{}}
  \toprule
  \thead{Purpose} & \thead{In-app} & \thead{Push (if the device opted in)} & \thead{Email} & \thead{Calendar}\\
  \midrule
  \endhead
  Something to keep (booked, moved, cancelled, marks released) & yes & yes & yes & the booking emails carry an \code{.ics} invite (same UID, rising \code{SEQUENCE}; \code{METHOD:CANCEL} on cancel), so the calendar entry updates itself and no separate ``add to calendar'' step is needed\\
  24 h reminder & yes & --- & yes (opt-out) & the calendar's own alert covers it\\
  1 h reminder & yes & yes & only when the student has no push device & ---\\
  Staff housekeeping (requests, returns) & yes & yes & daily digest instead of one email each & ---\\
  \bottomrule
\end{xltabular}}

Needs: outbox emails with attachments, and an opt-in for the invites. Until then \ui{Add to calendar} (Google, Outlook, \code{.ics}) is offered on each booking and calendar sync shows as \ui{coming soon}. Since \#142 the downloaded \code{.ics} carries alerts 1\,h and 10\,min before, the only reminder that fires offline with the browser closed.

Shipped from this list (ADR-018): no in-app echoes of a student's own booking actions, one delayed
confirmation email for the final state, superseded and coalesced items, and a bell that counts only
items that need the person. Still to do from the QA report of 18 Nov 2026:
\begin{itemize}
  \item \textbf{N1-digest} An email digest (daily, or one per assignment) and per-type email
    settings on the account page.
  \item \textbf{N1-email} Check real email volume and content for one demo in Mailpit (not covered
    by the QA run).
\end{itemize}

\section{Coming soon (shown disabled in the UI)}

\begin{itemize}
  \item Calendar sync (above).
  \item Google Meet links created automatically for online venues.
  \item More course timezones (only Pakistan time today).
\end{itemize}

\section{Open issues}

\begin{itemize}
  \item \textbf{Intermittent e2e timeout.} Once, the instructor test timed out on a page load after the service worker arrived (it passed on rerun). Watch for it; if it recurs, suspect \code{waitForLoadState("networkidle")} in \src{tests/e2e/helpers.ts} together with SW registration.
  \item \textbf{E2E now runs on a production build} (\code{next build && next start}, see \code{playwright.config.ts}). Under \code{next dev}, Fast Refresh from on-demand compiling sometimes dropped the page refresh after an action in multi-browser tests. \route{/dev/mail} is enabled there through \code{ENABLE_DEV_MAIL=1}.
\end{itemize}

\section[Hosting (AWS, \texttt{infra/aws})]{Hosting (AWS, \code{infra/aws})}

Bootstrap tier, about \$44/month against a \$55 budget (ADR-030). Next steps, in the order they're worth paying for:
\begin{itemize}
  \item \textbf{WAF} (\textasciitilde{}\$10/mo) or CloudFront's flat-rate plan once traffic is real; today API Gateway throttling is the only rate limit.
  \item \textbf{Multi-AZ RDS} (\code{db_multi_az}, \textasciitilde{}+\$14/mo) when a few minutes of downtime during an AZ outage stops being acceptable.
  \item \textbf{ALB instead of API Gateway} if responses need streaming, longer than 30 s, or larger than 10 MB.
\end{itemize}

